\section{Appendix: Liturgical Resolution}
This study concerns the Thirty-second Sunday in Ordinary Time, Year A, in the United States national calendar on 8 November 2026: an ordinary Sunday in green, with Gloria and Creed, Lectionary no. 154. The governing books are the \textit{Roman Missal, Third Edition, for Use in the Dioceses of the United States of America} (2011), based on the Latin third typical edition's emended reprint (2008), and the U.S. \textit{Lectionary for Mass}, second typical edition, Volume I. The permanent identity is PC-S58-A. The national calendar and the official dated readings establish this occurrence; no particular church's calendar has been supplied.

The common Missal prayers and antiphons are owned by the canonical Week 32 source record. The Year A leaf owns the Lectionary appointments and their branch resolution. The Sunday cycle is A; the independently resolved adjacent weekday cycle II supplies no reading to this Sunday. The General Introduction to the Lectionary, nos. 103--107, describes the Ordinary Time Old Testament selection in relation to the Gospel and the semi-continuous apostolic and Gospel courses. The Week 32 prayers are shared across cycles. Connections developed across the complete Mass are source-grounded theological synthesis; the cited Fathers expound their own biblical passages.

\begin{branchtable}
\bid{second-reading-long}; \bid{second-reading-short} & Lectionary 154; GIRM 360 & Appointed alternative & Second reading & 4:13--18 or 4:13--14; neither locally selected.\\
\bid{communion-psalm}; \bid{communion-luke} & Week 32; Ordinary Time rubric 6 & Appointed alternative & Communion & Both treated; neither assigned exclusively to Year A or asserted as locally selected.\\
\bid{entrance-chant-choice}; \bid{psalm-sung-choice}; \bid{acclamation-chant-choice} & U.S. GIRM 48, 61--63 & Permitted & Chants & No local replacement or omission established.\\
\bid{sprinkling}; \bid{preparation-chant-choice} & U.S. GIRM 51, 74 & Permitted & Entrance rite; preparation & No additional text selected.\\
\bid{preface-eucharistic-prayer-choice} & Ordinary Time rubric 5; GIRM 365 & Permitted & Preface and Eucharistic Prayer & Unselected; Prayer IV retains its integral Preface.\\
\bid{communion-chant-choice}; \bid{concluding-blessing-choice} & U.S. GIRM 87, 90 & Permitted & Communion; conclusion & No particular replacement chant or concluding blessing selected.\\
\end{branchtable}

No Sequence, short Gospel, proper Offertory antiphon, Sunday-specific Eucharistic Prayer insert or special ritual substitution is established for this target. Ordinary choices remain distinct from the ten text-bearing units inventoried above.

\section{Appendix: Scope and Qualifications}
The study develops two interpretations of the complete Year A appointment, including both second-reading lengths and both Communion antiphons. Its argument uses the passage contexts and the bounded reception research recorded in the owning dossier. It is an aid to Catholic study, not an official liturgical book or an authorized text for proclamation. No ecclesiastical approval is claimed.

Appointed-text completeness is rights-limited. Biblical English is the public-domain Douay--Rheims, Challoner revision, held through Gutenberg 1581; this electronic witness combines historical printings and is not identified as the 1899 American edition. The approved U.S. Lectionary controls proclamation. Biblical source passages beside adapted antiphons do not reproduce those adaptations. The approved response and acclamation, approved English orations, and full modern Latin orations are omitted; identifiers and original analytical commentary take their place. No new translation is offered for recitation.

The exact 2008 altar-book authentication remains incomplete. The shared Missal source owner records the secondary Latin reproduction, the official ICEL Antiphonary and FDLC excerpts, and a U.S. 2011 digital-book candidate whose publisher transmission has not been independently authenticated. The Prayer over the Offerings is discussed from the inspected Latin and that qualified candidate, without representing the candidate as an authorized publishable English exemplar. The older Collect conclusion in the retained English excerpts is not a current recitation text. Source-specific limitations remain with the common owner and are not resolved by this study's commentary.

Augustine and Chrysostom are read in the registered NPNF/CCEL historical English witnesses. Gregory is read in the identified Latin Wikisource witness; no claim of facsimile collation is made for those homilies. His arguments are paraphrased, not presented as a new liturgical translation. Bonaventure's Wisdom commentary was checked by research against the Quaracchi images. Aquinas is used through the identified Latin Corpus Thomisticum lectures in original analytical paraphrase. Bellarmine is used through O'Sullivan's historical abridgment, with its omissions acknowledged. A Father's quotation of another author does not establish that the quoted author's underlying work was independently inspected.

The reception search was bounded by the works and loci in the research dossier. It found no checked direct Father or saint exposition of the modern offerings and after-Communion wording. The treatments of those prayers are original analysis. The interpretations do not assign the oil uniquely to a sacrament or pronounce the fate of particular dead persons. Material differences concerning the sellers, inward purpose, nocturnal arrival and grief are retained at their claims.

This source edition contains the expansive study. Its concise study and homily are separate presentations under the same research owner. No print export or physical-page approval is claimed here; the sources and generated chronology are the basis for a selected export.

\section{References}
\begin{enumerate}
\item \textit{Douay--Rheims Bible}, Challoner revision, Gutenberg 1581 witness: Wisdom 6:13--17 (appointment 6:12--16); Psalms 22, 62 and 87 Vulgate; Matthew 24--25; Luke 24; 1 Thessalonians 4--5. Historical verse numbering and the Thessalonians offset are stated at the printed passages. Public-domain biblical text is excluded from the project's own content license.
\item United States Conference of Catholic Bishops, \href{https://bible.usccb.org/bible/readings/110826.cfm}{Thirty-second Sunday in Ordinary Time, 8 November 2026}, Lectionary 154; \href{https://www.usccb.org/resources/2026cal.pdf}{2026 Liturgical Calendar}, printed pp. 5, 43. Appointment and national occurrence control; approved English bodies are not reproduced.
\item \textit{Ordo lectionum Missae}, second typical edition (1981), no. 154, printed p. 85; \textit{General Introduction to the Lectionary}, nos. 103--107. Used for appointments and Ordinary Time relationships.
\item Common Missal source locators: \href{https://media.musicasacra.com/books/latin_missal2002.pdf}{Latin third-typical-edition digital reproduction}, Week XXXII, physical pp. 294--295, and Ordinary Time rubrics 2--6, physical p. 268; \href{https://www.liturgyoffice.org.uk/Missal/Music/Antiphonary.pdf}{ICEL Antiphonary}, printed p. 82, the entrance and both Communion alternatives; FDLC, \href{https://fdlc.org/wp-content/uploads/2026/05/ComboColPaC78e_ORDINARY_31-CTK_SPI2016.pdf}{\textit{Mystagogical Reflections}, Ordinary Time 31--Christ the King}, pp. 8--9, Collect and Prayer after Communion; \href{https://archive.org/download/Roman-Missal-Third-Edition/Roman\%20Missal.pdf}{U.S. 2011 digital-book candidate}, printed p. 492, Week 32, including the Prayer over the Offerings. The canonical Week 32 owner records the evidence and limitations; these locators supply no redistribution permission.
\item U.S. \textit{General Instruction of the Roman Missal}, current online state consulted 5 November 2026, \href{https://www.usccb.org/prayer-and-worship/the-mass/general-instruction-of-the-roman-missal/girm-chapter-2}{chapter II}, nos. 48, 51, 54, 61--64, 74, 87, 90, and \href{https://www.usccb.org/prayer-and-worship/the-mass/general-instruction-of-the-roman-missal/girm-chapter-7}{chapter VII}, nos. 354, 357, 360, 365. Used for the unselected option classes and current Collect conclusion.
\item Augustine, Sermon 93, §§1--17 (Benedictine numbering); Sermon 43 in \textit{Nicene and Post-Nicene Fathers}, first series, vol. 6, registered CCEL witness. The virginity, lamps, oil, sleep, sellers and final recognition argument.
\item Augustine, \textit{Expositions on the Psalms}, Psalms 22.1--6, 62.1--12 and 87.1--4; NPNF1, vol. 8, headings XXIII, LXIII and LXXXVIII, registered CCEL witness. Historical English translation; editorial annotations are distinct from Augustine's text.
\item John Chrysostom, \textit{Homilies on Matthew}, 77.1--3 and 78.1--2, NPNF1, vol. 10, registered CCEL witness; \textit{Homilies on First Thessalonians}, 6--8, opening pericope expositions only, NPNF1, vol. 13. Historical English translations.
\item Gregory the Great, \textit{Homiliae in Evangelia}, 12.1--6 and 23.1--2, \href{https://la.wikisource.org/wiki/Homiliarum_in_Evangelia}{Latin Wikisource witness}. The retained export credits Gregorius Magnus and contributor Mizardellorsa and carries CC BY-SA 3.0 for its transcription/encoding; the ancient Latin is a distinct rights object. This study supplies original analytical paraphrase.
\item Bonaventure, \textit{Commentarius in librum Sapientiae}, chapter 6, Vulgate vv. 13--17, \textit{Opera omnia}, Quaracchi VI (1893), pp. 147--149. Medieval exegesis read in the historical facsimile.
\item Thomas Aquinas, \textit{Super I ad Thessalonicenses}, chapter 4, lectio 2; \textit{Super Matthaeum}, chapter 25, lectio 1 and opening of lectio 2, identified Latin Corpus Thomisticum witnesses. The modern site bodies are not reproduced.
\item Robert Bellarmine, \textit{A Commentary on the Book of Psalms}, John O'Sullivan's historical English abridgment, Psalms 22 and 62:1--8. The translator's stated omissions qualify the extent of this witness.
\end{enumerate}
